\section{Evaluation and validity}
\label{sec:methodology}

\subsection{Experimental unit and controlled variables}

For the primary study, the experimental unit is one completed inference
run on one allocated node under one memory policy.  The HBM and DDR
treatments must use the same:

\begin{itemize}
  \item physical socket, active cores and thread affinity;
  \item model file, tokenizer, quantization and model hash;
  \item runtime commit, compiler, libraries, build flags and AMX path;
  \item prompt tokens, generation length, batch/concurrency and decoding
        settings;
  \item page size, memory mapping/locking policy and warm-up protocol;
        and
  \item metric definitions and measurement code.
\end{itemize}

Only the memory-node policy changes.  If HBM and DDR reside under
different CPU affinities, the correct comparison pairs each memory node
with the same local cores or deliberately fixes cores and documents the
distance.  The platform topology, not a node-number naming convention,
determines the pair.

\subsection{Placement validation}

Placement is a measured intermediate outcome.  A run is included in the
causal HBM/DDR comparison only when its important pages reside on the
intended node and remote traffic remains within a pre-specified
tolerance.

A practical validation protocol is:

\begin{enumerate}
  \item record \texttt{lscpu}, \texttt{numactl --hardware}, NUMA
        distances and available memory immediately before the run;
  \item launch with explicit CPU and memory binding;
  \item capture page counts per NUMA node after model loading, after
        warm-up and during the measured phase;
  \item record whether weights are mapped or copied, and whether packed
        representations create additional allocations;
  \item log automatic NUMA balancing, transparent huge pages and any
        migration observed; and
  \item retain failed placement checks as failed trials rather than
        silently replacing their values.
\end{enumerate}

System-wide page-cache eviction is disruptive on a shared cluster and
must not be assumed available.  Alternatives include exclusive nodes,
administrator-approved cache handling, a controlled non-mmap load path,
or distinct copied model files.  Whichever protocol is chosen must be
tested for equivalent numerical output and documented.  Reversing the
treatment order and interleaving HBM/DDR trials helps detect residual
cache and thermal drift.

\subsection{Workload phases and metrics}

Prefill and decode should be delimited by the runtime or by a validated
instrumentation layer.  The primary metrics are:

\begin{description}[leftmargin=3.1cm, style=nextline, itemsep=4pt]
  \item[TTFT] wall-clock time from request admission to the first
        generated token, with tokenization and model load explicitly
        included or excluded;
  \item[prefill throughput] prompt tokens processed per second during
        the prefill interval;
  \item[ITL] time between successive generated tokens after the first,
        reported as a distribution and not only a mean;
  \item[decode throughput] generated tokens per second for the measured
        steady interval; and
  \item[memory evidence] achieved HBM/DDR bandwidth, remote traffic,
        peak resident memory and per-node page counts.
\end{description}

If a single run contains both phases, its end-to-end latency may be
reported as a user metric, but it cannot replace phase measurements.
Early decode tokens may behave differently from later tokens as the KV
cache grows; the analysis should state which tokens are excluded as
warm-up and which define the sustained interval.

\subsection{Experiment matrix and repetition}

Pilot runs should determine the feasible matrix rather than fixing many
arbitrary combinations in advance.  The final design should nevertheless
contain:

\begin{itemize}
  \item at least one configuration comfortably below HBM capacity;
  \item at least one second model or quantization to test whether the
        effect follows live-set size rather than a single model;
  \item two or more arithmetic-intensity points, for example sequence
        concurrency or batch; and
  \item context/capacity points concentrated near any observed
        transition, with fewer points far from it.
\end{itemize}

Runs are performed in blocks on the same node allocation where
possible.  Within a block, treatment order is randomised or balanced.
The number of repetitions is set from pilot variance and a target
precision for the HBM/DDR effect, not from a fixed convention such as
three runs.

For each configuration, report all valid observations, median and a
dispersion measure, plus a confidence interval for the paired or
blocked treatment effect.  Means may be added when the distribution is
well behaved.  Outlier rules must be defined from operational evidence
(for example, a logged placement failure or node interruption), not
from an undesirable performance value.

\subsection{Mechanism analysis}

A credible explanation triangulates three measurements:

\begin{enumerate}
  \item \textbf{placement:} the intended pages are local to the selected
        tier;
  \item \textbf{pressure:} the measured phase issues enough traffic to
        make memory bandwidth a plausible limit; and
  \item \textbf{response:} the user metric changes consistently with
        the achieved resource.
\end{enumerate}

The arithmetic intensity estimate
\[
  I_{\mathrm{observed}} =
  \frac{\text{estimated or counted operations}}
       {\text{measured memory bytes}}
\]
can be compared with a machine balance derived from local
microbenchmarks.  Because operations are difficult to count exactly in
quantized kernels, the thesis should state whether the numerator comes
from model arithmetic, hardware events or runtime instrumentation.
The conclusion should not rely on a precise Roofline position when the
counter mapping is uncertain.

The HBM/DDR speedup is
\[
  S_{\mathrm{tier}} =
  \frac{T_{\mathrm{DDR}}}{T_{\mathrm{HBM}}}
\]
for latency, or the inverse ratio for throughput.  Compare it with the
measured bandwidth ratio as a bound and explanatory reference, not as
a target.  If the speedup changes with batch or context, show whether
achieved bandwidth, compute utilisation or live set changes at the same
point.

\subsection{Correctness and output equivalence}

Memory placement should not change model semantics.  Identical seeds,
deterministic decoding where supported, token sequences and numerical
checks are required across HBM and DDR.  A runtime patch additionally
needs:

\begin{itemize}
  \item allocation tests for requested and insufficient tier capacity;
  \item page-residency tests for each buffer type;
  \item equality or tolerance checks against the unmodified runtime;
  \item clear logging of the selected policy and actual allocation; and
  \item failure rather than silent fallback when an experiment requests
        strict placement.
\end{itemize}

Quality differences between quantizations are not caused by HBM, but
they matter when comparing model formats.  Cross-quantization
performance results should therefore be labelled as separate operating
points and accompanied by an appropriate quality check or a citation
to the model's validated representation.

\subsection{Threats to validity}

\paragraph{Internal validity.}
First-touch, file-backed pages, migration, remote memory, frequency
changes and different kernel paths can confound the tier comparison.
The placement and build checks above address these threats.  The study
must avoid comparing HBM and DDR on different sockets unless socket is
modelled as a block.

\paragraph{Construct validity.}
Tokens per second alone does not represent interactive inference.
TTFT and ITL are necessary.  STREAM bandwidth is not LLM bandwidth;
hardware counters and phase timing provide the connecting evidence.
Nominal model size is not resident live set.

\paragraph{Statistical conclusion validity.}
Autocorrelation within an allocation, node-to-node variation and
thermal drift make individual runs non-independent.  Blocking by node
and allocation, balanced order and confidence intervals over the
appropriate unit reduce false precision.

\paragraph{External validity.}
One CPU generation, one runtime and a limited model set do not establish
all CPU or LLM behaviour.  Results should be stated for the tested
Max~9480 configurations and explained through measurable variables so
that later systems can test transfer.

\paragraph{Software validity.}
Upstream \lcpp{} changes quickly.  The primary dataset must use a frozen
commit and archived build manifest.  A later version can be a
replication point.  Static source observations must not be presented as
runtime behaviour without the validation described in
\secref{sec:audit}.

\subsection{Falsification criteria}

The programme is designed to survive negative evidence.  The following
outcomes revise its interpretation:

\begin{itemize}
  \item If valid HBM and DDR runs have equivalent phase performance and
        neither approaches memory bandwidth, the workload point is not
        memory-bound; the thesis should report the boundary rather than
        claim a software failure.
  \item If they are equivalent while both show high bandwidth pressure,
        investigate cache behaviour, counters and hidden placement
        before accepting the null effect.
  \item If DDR is faster under light load, this is consistent with the
        latency evidence in Ibeid et al.\ and should be explained rather
        than discarded.
  \item If unmanaged two-socket execution loses but local replicas
        scale aggregate throughput, the issue is placement and workload
        decomposition, not a universal single-socket advantage.
  \item If existing operating-system controls express the best selective
        policy, a runtime patch has no thesis justification.
  \item If a patch changes placement correctly but not performance, the
        software capability exists but the proposed policy lacks a
        performance contribution for the tested regime.
\end{itemize}

\subsection{Immediate actions}

The next work should be empirical preparation, not additional general
background prose:

\begin{enumerate}
  \item confirm access to a Max~9480 node and the available Flat/Cache
        and Quadrant/SNC settings with CRESCO8 administrators;
  \item capture one complete topology and memory-capacity manifest;
  \item build the pinned \lcpp{} snapshot with reproducible compiler
        flags and verify AMX features;
  \item test whole-process HBM and DDR placement with a small model and
        inspect page residency;
  \item measure local and deliberately remote memory bandwidth under
        the same core binding; and
  \item use pilot variance and observed usable HBM to freeze the Stage~1
        model, quantization, contexts and repetitions.
\end{enumerate}

After these checks, a new dated supervisor briefing in
\texttt{research/briefings/} can be written with an agreed primary
research question and a feasible experiment matrix.  Detailed commands,
raw manifests and exploratory results belong in a reproducibility
repository or thesis appendix rather than in the literature-review
chapter.
